\begin{lemma}[A H\"older datum at the fixed exponent $\tfrac12$, uniform in the centre]
Let $n\ge2$, $p>1$, $0<\tau\le p/2$, $d>0$, $M\ge0$ and $N_2\ge0$, and write
\[
q=\frac n{n-1}\ (>1),\qquad
\lambda^\natural=n+\Bigl(1-\frac\tau p\Bigr)q .
\]
Assume the cited Campanato embedding (C5) at the exponent $q$ on the single ball
$B_{d/4}(0)$ (\noderef{CitedCampanatoEmbedding}).  Then there is a constant
$C_1\ge0$, depending only on $n,p,\tau,d,M$ and $N_2$, with the following
property.

Let $\Omega\subseteq\R^n$ be open, let $\Omega'\subseteq\R^n$ be a set with
$d\le\operatorname{dist}(x,y)$ for every $x\in\Omega'$ and every
$y\notin\Omega$, and let $u,U\colon\R^n\to\R$ satisfy
\begin{itemize}
\item[(a)] $|u|\le M$ almost everywhere on $\Omega$;
\item[(b)] $U=u$ almost everywhere on $\Omega$;
\item[(c)] $U$ is continuous on $B_{d/4}(y_0)$ for every $y_0\in\Omega'$;
\item[(d)] $\|u\|_{\mathcal L^{q,\lambda^\natural}(B_{d/4}(y_0))}\le N_2$ for
  every $y_0\in\Omega'$ (\noderef{CampanatoBoundOn}).
\end{itemize}
Then for every $x_0\in\Omega'$ and every $R$ with $0<R\le d/16$ one has
$U\in C^{1/2}\bigl(\overline B_R(x_0)\bigr)$ (\noderef{MemHolderOn}) and
\[
|U(x)-U(y)|\le C_1\operatorname{dist}(x,y)^{1/2}
\qquad\bigl(x,y\in\overline B_R(x_0)\bigr),
\]
with the \emph{same} $C_1$ for every centre $x_0$ and every such radius $R$.
\end{lemma}
\begin{proof}
Throughout, $\omega_n=|B_1(0)|\in(0,\infty)$ is the volume of the unit ball
(Preamble, the volume of a ball), so that $|B_\rho(z)|=\omega_n\rho^{\,n}$ for
every centre $z$ and every $\rho\ge0$.

\emph{Step 1: the constants, chosen before the data.}  Put
\[
s=\frac{\lambda^\natural-n}q=1-\frac\tau p .
\]
Since $0<\tau\le p/2$ and $p>1$ we have $0<\tau/p\le\tfrac12$, hence
\[
\tfrac12\le s<1 ,
\]
and in particular $s-\tfrac12\ge0$.  Also $n\ge2$ gives $n-1\ge1>0$, so
$q=n/(n-1)>1$ is well defined and positive.

The ball $B_{d/4}(0)$ is a bounded Lipschitz domain
(\noderef{IsBoundedLipschitzDomain}): this is \noderef{BallIsLipschitzDomain},
applied with the present $n$ --- its hypothesis $1\le n$ follows from $2\le n$
--- at the centre $0$ and the radius $d/4$, whose positivity follows from
$d>0$.  That discharges the Lipschitz antecedent with which the assumed
predicate $\mathrm{CitedCampanatoEmbedding}\;n\;q\;B_{d/4}(0)$
(\noderef{CitedCampanatoEmbedding}) opens.  Read the predicate at the exponent
$\lambda^\natural$ and at the constant
\[
N^\natural=\max\Bigl\{N_2,\ M\bigl(\omega_n(d/4)^n\bigr)^{1/q}\Bigr\}\ \ge0 .
\]
Its three guards hold: $n<\lambda^\natural$, because
$\bigl(1-\tfrac\tau p\bigr)q=sq>0$ as $s\ge\tfrac12>0$ and $q>0$;
$\lambda^\natural<n+q$, because $sq<q$ as $s<1$ and $q>0$; and
$0\le N^\natural$, because $N_2\ge0$.  It produces a constant $N_h\ge0$, and
that constant is a function of $n,q,d,\lambda^\natural$ and $N^\natural$ alone,
hence of $n,p,\tau,d,M$ and $N_2$ alone, because $q$ is a function of $n$,
$\lambda^\natural$ a function of $n,p,\tau$, $\omega_n$ a function of $n$ and
$N^\natural$ a function of $n,q,d,M$ and $N_2$.  Put
\[
C_1=N_h\max\{1,d\}^{\,s-\frac12}\ \ge0 ,
\]
which is likewise a function of $n,p,\tau,d,M$ and $N_2$ alone, and which is
nonnegative because $N_h\ge0$ and $\max\{1,d\}^{\,s-1/2}>0$.

\emph{No further choice is made below}: everything that follows is proved for
the data, with this one $C_1$.  Fix now $\Omega,\Omega',u,U$ as in the statement,
and fix a point $y_0\in\Omega'$.  Write $B^{y_0}=B_{d/4}(y_0)$.

\emph{Step 2: the ball $\overline{B^{y_0}}$ lies in $\Omega$.}  Let
$z\in\overline B_{d/4}(y_0)$, so $\operatorname{dist}(z,y_0)\le d/4$, and
suppose $z\notin\Omega$.  Then the separation hypothesis, applied at
$x=y_0\in\Omega'$ and $y=z\notin\Omega$, gives
$d\le\operatorname{dist}(y_0,z)\le d/4$, hence $d\le d/4$ and so
$\tfrac34d\le0$, contradicting $d>0$.  Therefore
$\overline B_{d/4}(y_0)\subseteq\Omega$, and a fortiori
$B^{y_0}\subseteq\Omega$.

\emph{Step 3: the translate, and the two function hypotheses of \textup{(C5)} at
the one constant $N^\natural$.}  Write $u_{y_0}(z)=u(z+y_0)$.  Lebesgue measure
and the metric on $\R^n$ are invariant under the translation
$z\mapsto z+y_0$ (Preamble, invariance of Lebesgue measure under translations;
and $\|(z+y_0)-(z'+y_0)\|=\|z-z'\|$), and
\[
B_{d/4}(0)+y_0=B^{y_0},\qquad B_\rho(x)+y_0=B_\rho(x+y_0)
\qquad(x\in\R^n,\ \rho>0),
\]
because $\|z-x\|<\rho$ iff $\|(z+y_0)-(x+y_0)\|<\rho$.  Consequently, for every
measurable $S\subseteq B_{d/4}(0)$ the Preamble identities $|S+y_0|=|S|$ and
$\int_{S+y_0}w=\int_Sw(\cdot+y_0)$ carry every integral, every measure and
every mean value of $u$ over $B^{y_0}\cap B_\rho(x+y_0)$ to the corresponding
one of $u_{y_0}$ over $B_{d/4}(0)\cap B_\rho(x)$, and the invariance of the
metric gives
$\operatorname{diam}B_{d/4}(0)=\operatorname{diam}B^{y_0}$.  Also
$|u_{y_0}|\le M$ almost everywhere on $B_{d/4}(0)$: the set
$\{z\in B_{d/4}(0):|u(z+y_0)|>M\}$ is the translate by $-y_0$ of
$\{z\in B^{y_0}:|u(z)|>M\}$, which is contained in
$\{z\in\Omega:|u(z)|>M\}$ by Step 2 and is therefore null by hypothesis (a); and
a translate of a null set is null by the same Preamble item.  Hence:
\begin{itemize}
\item $\mathrm{CampanatoBoundOn}\;n\;q\;\lambda^\natural\;B_{d/4}(0)\;u_{y_0}\;
  N^\natural$ (\noderef{CampanatoBoundOn}).  Its $L^q$ clause transports
  directly from hypothesis (d) at the centre $y_0$, being $N$-free: $u\in
  L^q(B^{y_0})$ gives $u_{y_0}\in L^q(B_{d/4}(0))$ by the Preamble's
  translation identity for $L^q$ seminorms.  Its inequality clause: let
  $x\in B_{d/4}(0)$ and $0<\rho<\operatorname{diam}B_{d/4}(0)$.  Then
  $x+y_0\in B^{y_0}$ and $\rho<\operatorname{diam}B^{y_0}$, so hypothesis (d) at
  the centre $y_0$ applies at the point $x+y_0$ and the radius $\rho$ and gives
  \[
  \int_{B^{y_0}\cap B_\rho(x+y_0)}\Bigl|u-\frac1{|B^{y_0}\cap
  B_\rho(x+y_0)|}\int_{B^{y_0}\cap B_\rho(x+y_0)}u\Bigr|^q
  \le N_2^{\,q}\rho^{\,\lambda^\natural},
  \]
  and the left-hand side equals the corresponding integral for $u_{y_0}$ over
  $B_{d/4}(0)\cap B_\rho(x)$ by the transport just described.  Finally
  $N_2^{\,q}\rho^{\lambda^\natural}\le(N^\natural)^q\rho^{\lambda^\natural}$,
  because $0\le N_2\le N^\natural$, $q>0$ and $\rho^{\lambda^\natural}\ge0$.
\item $\mathrm{MorreyBoundOn}\;n\;q\;0\;B_{d/4}(0)\;u_{y_0}\;N^\natural$
  (\noderef{MorreyBoundOn}).  Both clauses come from
  $|u_{y_0}|\le M$ a.e.\ on $B_{d/4}(0)$.  The $L^q$ clause holds because a
  function bounded almost everywhere on a set of finite measure lies in $L^q$ of
  that set, and $|B_{d/4}(0)|=\omega_n(d/4)^n<\infty$.  For the inequality
  clause let $x\in B_{d/4}(0)$ and $0<\rho<\operatorname{diam}B_{d/4}(0)$; then
  \[
  \int_{B_{d/4}(0)\cap B_\rho(x)}|u_{y_0}|^q
  \le M^q\bigl|B_{d/4}(0)\bigr|=M^q\omega_n(d/4)^n
  \le(N^\natural)^q=(N^\natural)^q\rho^{\,0},
  \]
  the first step by monotonicity of the integral over
  $B_{d/4}(0)\cap B_\rho(x)\subseteq B_{d/4}(0)$ for the nonnegative integrand
  $|u_{y_0}|^q\le M^q$ a.e., the third because
  $N^\natural\ge M(\omega_n(d/4)^n)^{1/q}\ge0$ and $t\mapsto t^q$ is
  nondecreasing on $[0,\infty)$, and the last because $\rho^{\,0}=1$ for
  $\rho>0$.
\end{itemize}

\emph{Step 4: what \textup{(C5)} returns, and its transport back.}  All three
terms of its conclusion are used.  The predicate binds $N_h$ \emph{before} the
function it is read at, so the $N_h$ of Step 1 --- drawn once, at the ball
$B_{d/4}(0)$, the exponent $\lambda^\natural$ and the constant $N^\natural$, all
three of which are independent of $y_0$ --- serves every $y_0\in\Omega'$.
Applied to the function $u_{y_0}$, for which the two hypotheses were just
verified at the constant $N^\natural$, it gives a function $w^0$ with
\begin{itemize}
\item $u_{y_0}=w^0$ almost everywhere on $B_{d/4}(0)$;
\item $w^0\in C^{s}\bigl(\overline{B_{d/4}(0)}\bigr)$ (\noderef{MemHolderOn}),
  the exponent being $(\lambda^\natural-n)/q=s$, so in particular $w^0$ is
  continuous on the closure of $B_{d/4}(0)$; and
\item $|w^0(x)-w^0(y)|\le N_h\operatorname{dist}(x,y)^{s}$ for all
  $x,y$ in the closure of $B_{d/4}(0)$, with the $N_h$ of Step 1.
\end{itemize}
Put $w_{y_0}(z)=w^0(z-y_0)$.  Then $u=w_{y_0}$ almost everywhere on $B^{y_0}$,
because $\{z\in B^{y_0}:u(z)\ne w_{y_0}(z)\}$ is the translate by $y_0$ of
$\{z\in B_{d/4}(0):u_{y_0}(z)\ne w^0(z)\}$ and a translate of a null set is
null.  Moreover $z\mapsto z-y_0$ is a homeomorphism of $\R^n$ carrying
$B^{y_0}$ onto $B_{d/4}(0)$, hence carrying the closure of $B^{y_0}$ into the
closure of $B_{d/4}(0)$, so $w_{y_0}$ is continuous on the closure of
$B^{y_0}$; and, since that map is an isometry,
\[
|w_{y_0}(x)-w_{y_0}(y)|\le N_h\operatorname{dist}(x,y)^{s}
\]
for all $x,y$ in the closure of $B^{y_0}$.

\emph{Step 5: $U=w_{y_0}$ pointwise on $B^{y_0}$.}  By hypothesis (c), $U$ is
continuous on the open set $B^{y_0}$; and $w_{y_0}$ is continuous on the closure
of $B^{y_0}$, hence on $B^{y_0}$.  By hypothesis (b) and Step 2, $U=u$ almost
everywhere on $B^{y_0}$, and by Step 4, $u=w_{y_0}$ almost everywhere there; so
$U=w_{y_0}$ almost everywhere on $B^{y_0}$.  The set
\[
E=\{x\in B^{y_0}:U(x)\ne w_{y_0}(x)\}
\]
is the preimage of $\R\setminus\{0\}$ under the function $U-w_{y_0}$, which is
continuous on the open set $B^{y_0}$, so $E$ is open in $B^{y_0}$ and hence open
in $\R^n$; and $E$ is Lebesgue-null by the previous sentence.  If $E$ were
nonempty it would contain some $B_\rho(z)$ with $\rho>0$, whose measure
$\omega_n\rho^{\,n}$ is positive, contradicting $|E|=0$.  So $E=\emptyset$, i.e.
\[
U=w_{y_0}\quad\text{pointwise on }B^{y_0}.
\]

\emph{Step 6: the conclusion, at the fixed exponent $\tfrac12$.}  Let
$x_0\in\Omega'$ and $0<R\le d/16$, and apply Steps 2--5 at $y_0=x_0$.  Since
$R\le d/16<d/4$ we have
\[
\overline B_R(x_0)\subseteq B_{d/4}(x_0)=B^{x_0},
\]
so by Step 5, $U=w_{x_0}$ pointwise on $\overline B_R(x_0)$, and
$\overline B_R(x_0)$ is contained in the closure of $B^{x_0}$.  Two consequences.

First, $U$ is continuous on $\overline B_R(x_0)$: by hypothesis (c) at the
centre $x_0$, $U$ is continuous on $B^{x_0}$, and continuity on a set restricts
to every subset of it.

Second, let $x,y\in\overline B_R(x_0)$ and put
$t=\operatorname{dist}(x,y)$.  If $t=0$ then $x=y$ and both sides of the
asserted inequality are $0$, the right-hand side because $0^{1/2}=0$ under the
real-power convention the Preamble records ($0^{r}=0$ for $r\ne0$).  If $t>0$
then, $x$ and $y$ lying in the closure of $B^{x_0}$, Step 4 gives
$|U(x)-U(y)|=|w_{x_0}(x)-w_{x_0}(y)|\le N_ht^{\,s}$; and
\[
t=\operatorname{dist}(x,y)\le\operatorname{dist}(x,x_0)+
\operatorname{dist}(x_0,y)\le2R\le\frac d8\le d\le\max\{1,d\},
\]
so, $r\mapsto r^{\,s-1/2}$ being nondecreasing on $(0,\infty)$ for the
nonnegative exponent $s-\tfrac12$,
\[
t^{\,s}=t^{1/2}\,t^{\,s-\frac12}\le t^{1/2}\max\{1,d\}^{\,s-\frac12}.
\]
Multiplying by $N_h\ge0$,
\[
|U(x)-U(y)|\le N_h\max\{1,d\}^{\,s-\frac12}\,t^{1/2}
=C_1\operatorname{dist}(x,y)^{1/2},
\]
which is the displayed inequality.  Together with the continuity of $U$ on
$\overline B_R(x_0)$ and $C_1\ge0$, it is also exactly
$U\in C^{1/2}(\overline B_R(x_0))$ (\noderef{MemHolderOn}), whose two clauses
are continuity on the set and the existence of a nonnegative constant for which
the two-point bound holds.  The constant $C_1$ was fixed in Step 1, before
$\Omega,\Omega',u,U$ and before $x_0$ and $R$, so it is one constant for every
centre and every admissible radius, as asserted.
\end{proof}
